\section{Optimality: a finite window for an infinite problem}
\label{sec:optimality}

This chapter answers one question: in what sense is the finite-horizon problem
$\mathcal{P}_N(x_k)$ a substitute for the infinite-horizon problem \eqref{eq:infhor}? The answer
goes through the value function and the principle of optimality, and it identifies the terminal
cost $V_f$ as the ingredient that makes the substitution exact. Throughout the chapter three
assumptions hold: the model is known and perfect, the state is measured, there are no
disturbances.

\subsection{The value function and Bellman's principle}

\begin{definition}[value function]
The \emph{value function} $V(x)$ is the optimal cost of the infinite-horizon problem
\eqref{eq:infhor} started at $x$: the cost of controlling the system optimally from $x$ onwards.
It is $+\infty$ when no admissible input sequence exists from $x$.
\end{definition}

\begin{theorem}[principle of optimality]
Let $\{u^*_0, u^*_1, \dots\}$ be an optimal input sequence for the problem started at $x_0$, and
let $x^*_N$ be the state it reaches at time $N$. Then the tail $\{u^*_N, u^*_{N+1}, \dots\}$ is an
optimal input sequence for the problem started at $x^*_N$.
\end{theorem}

The reason is simple: if a cheaper tail existed from $x^*_N$, replacing the tail of the original
sequence with it would give a cheaper sequence from $x_0$, against the optimality of the
original. The consequence is a decomposition of the optimal cost at any junction time $N$:
\begin{equation}
  V(x_0) = \sum_{i=0}^{N-1} \ell(x^*_i, u^*_i) + V(x^*_N).
  \label{eq:bellman}
\end{equation}
The optimal cost from $x_0$ is the cost of the first $N$ steps plus the value function of the
state reached: the infinite tail is summarized by one number.

\subsection{Dual-mode control}

Consider the constrained linear-quadratic problem. Assume the system controllable and the origin
strictly inside $\X$ and $\U$. Then an optimal sequence exists and its cost is finite; a finite
sum of positive-definite terms forces $x_k \to 0$ and $u_k \to 0$, so the constraints can be
active only in an initial phase. After some finite time $N$ the state is close enough to the
origin that the trajectory generated by the unconstrained law $u = Kx$ respects all the
constraints from then on, and by the principle of optimality applied to the tail problem that
law is not only admissible there, it is \emph{optimal}, because from that point on the
constrained and the unconstrained problems have the same solution.

This is \emph{dual-mode control}: the input is free in the first $N$ steps, and it is the LQR
feedback afterwards (\cref{fig:dualmode}). The tail is then known in closed form, and so is its
cost, $x_N\T W x_N$ from \eqref{eq:Wcost}. Substituting it in \eqref{eq:bellman}, the infinite
problem becomes
\begin{equation}
  \min_{u_0,\dots,u_{N-1}} \ \sum_{i=0}^{N-1}\big(x_i\T Q x_i + u_i\T R u_i\big) + x_N\T W x_N,
  \label{eq:dualmode}
\end{equation}
subject to the dynamics and the constraints over the first $N$ steps. This problem has $N$
decision variables and it is solvable. Its solution coincides with the first $N$ inputs of the
infinite-horizon optimum, $u^{*N}_k = u^{*\infty}_k$ for $k = 0,\dots,N-1$, provided $N$ is
large enough for the tail to be unconstrained.

\begin{figure}[t]
\centering
\begin{tikzpicture}[x=1cm,y=1cm,>=Stealth]
  \draw[->] (0,0) -- (11,0) node[right] {time};
  \draw[->] (0,0) -- (0,2.4) node[above] {cost};
  % first phase: free moves
  \fill[mpc!25] (0,0) rectangle (4.5,1.6);
  \draw[thick,mpc] plot[smooth] coordinates {(0,2.0) (1,1.6) (2,1.25) (3,1.0) (4.5,0.8)};
  \node[mpc,font=\small] at (2.25,0.45) {$\sum_{i=0}^{N-1}\ell(x_i,u_i)$};
  \node[mpc,font=\small,above,align=center] at (2.25,2.05) {mode 1: free inputs,\\constraints may be active};
  % tail
  \fill[sun!25] (4.5,0) rectangle (10.5,1.6);
  \draw[thick,sun] plot[smooth] coordinates {(4.5,0.8) (6,0.5) (7.5,0.3) (9,0.15) (10.5,0.08)};
  \node[sun!70!black,font=\small] at (7.5,0.9) {$V(x_N) = x_N\T W x_N$};
  \node[sun!70!black,font=\small,above,align=center] at (7.5,2.05) {mode 2: $u = Kx$,\\constraints inactive};
  \draw[dashed] (4.5,0) -- (4.5,2.3);
  \node[below] at (4.5,0) {$N$};
  \node[below] at (0,0) {$0$};
\end{tikzpicture}
\caption{Dual-mode control. In the first $N$ steps the inputs are free and the constraints may
be active; from $N$ on the unconstrained LQR feedback is optimal, and the cost of that tail is
the quadratic form $x_N\T W x_N$, which becomes the terminal cost of a finite-horizon problem.}
\label{fig:dualmode}
\end{figure}

\subsection{The terminal cost as the cost-to-go}

Problem \eqref{eq:dualmode} is $\mathcal{P}_N(x_k)$ with the quadratic stage cost and with
terminal cost $V_f(x) = x\T W x$. The lesson of the construction is general: the terminal cost
of an MPC problem should be the \emph{value function} of the infinite-horizon problem, the cost
of what happens after the window. When it is, the finite problem and the infinite one have the
same solution, and solving the finite problem again at every sampling time --- which is MPC ---
reproduces the infinite-horizon optimum. This is the precise sense in which MPC is optimal.

The practical difficulty is that $V(x)$ is unknown in general. The unconstrained
linear-quadratic problem is the exception: there $V(x) = x\T W x$ exactly, and this is why, in
linear-quadratic MPC, the terminal cost is taken from the Riccati or the Lyapunov equation and
the horizon is taken long enough. In other problems $V(x)$ has to be approximated, or learned:
the value function is also the object computed by dynamic programming, and learning it from
experience is the subject of reinforcement learning, the second part of the course.

\begin{remark}[the terminal weight matters]
The Service Area Lab on the website shows what happens when the terminal cost is not the value
function. With a window of twenty minutes and a terminal weight equal to the stage weight on the
stored energy, the controller is myopic: it sees almost no value in the state it leaves at the
end of the window, and the storage drifts away from its reference. With the terminal weight
from the Riccati equation, about four hundred times larger, the same twenty-minute controller
behaves like the infinite-horizon one.
\end{remark}

\subsection{From one step to the next}

Under the nominal assumptions, re-solving the problem at time $k+1$ returns the plan the
controller already had: $u^*_{\ik{1}}$ becomes $u^*_{0|k+1}$, and so on. The controller
confirms itself, and the receding horizon changes nothing. In reality the model is not perfect,
the forecasts are wrong and disturbances act on the process, so the state measured at $k+1$ is
not the one predicted; this is why the plan is rebuilt at every step rather than computed once.

The cost actually paid by the controller in closed loop is the sum of the stage costs along
the trajectory it generates, $J^{cl}(x_0) = \sum_{k \ge 0} \ell(x_k, \kappa_N(x_k))$. Its
difference from $V(x_0)$ is the \emph{optimality gap}; it is zero in the ideal conditions of
this chapter, and it measures what is lost when they fail.

\begin{reading}
\refkey{Cannon\_2016}, section 2.3, ``The Dual-Mode Prediction Paradigm'' (pp.~18--21);
Theorem~2.2 there is the formal version of the result of this chapter. \refkey{Rawlings\_2017},
sections 1.3.3 and 1.3.4 (pp.~18--23) and section 2.3, ``Dynamic Programming Solution''
(pp.~107--112), for the same material reached through dynamic programming, which is the road
taken again in the reinforcement learning part. Studying the proofs is optional.
\end{reading}
